\scanpage{8}
% archive foliation: 7
Soit $X$ schéma lisse sur $S$, \deleted{Un} $n$ un entier $\geq0$. Une \uline{$n$-connexion projective} sur $X$ (rel. à $S$) est la donnée

 a) d'un fibré de B.S $P$ sur $X$ avec une section, ou ce qui revient au même, une extension
\diagram{\begin{tikzcd}[column sep=small]
0\arrow[r]&\Omega\arrow[r]&E\arrow[r]&\mathcal{O}_X\arrow[r]&0
\end{tikzcd}}
de Modules loc. libres sur $X$, ou encore un couple $(\Omega,\mathcal{T})$ de $\Omega$ loc. libre sur $X$ et $\mathcal{T}$ un $\Omega$-tors.

 b) Un \inserted{$X$-}isom des voisinages inf. d'ordre $n$ de la diagonale de $X\times_S X$ avec le vois. d'ordre $n$ de $\sigma$ dans $P$, ou ce qui revient au même, un isom de \inserted{$\mathcal{O}_X$-}Algèbres augmentées
\[
P^n_{X/S}\overset{\alpha_n}{\simeq}\mathcal{O}_{\widehat P}/\widehat{\mathcal{J}}^{n+1}
\quad(\simeq\mathcal{T}\times_{\Omega}\operatorname{Sym}^{\bullet}(\Omega)/\operatorname{Sym}^{+}(\Omega)^{n+1})
\]
(où $\widehat P$ est le complété formel de $P$ le long de $\sigma(X)$, et $\widehat{\mathcal{J}}$ l'Idéal de déf. de la section $\sigma$).

Pour $n=0$, la donnée b est vide. Pour $n=1$, elle revient à un iso
\[
(*)\qquad\Omega\simeq\Omega^1_{X/S}
\]
\margin{donc le rang de $\Omega$ en tt pt doit être la dim rel de $X$ sur $S$ en ce pt.}

Pour $n=2$, la structure de $P^2_{X/S}$ est entièrement déterminée par la structure d'extension de $\Omega^1_{X/S}$ par $\operatorname{Sym}^2\Omega^1_{X/S}$, \deleted{\unsure{qui doit|qui définit}} \inserted{i.e.} \deleted{\unsure{d'une connex}} un torseur sous $\underline{\operatorname{Hom}}_{\mathcal{O}_S}(\Omega^1_{X/S},\operatorname{Sym}^2\Omega^1_{X/S})$ lequel doit donc être muni d'un iso avec le torseur déduit de $\mathcal{T}$ sous $\Omega$, de
\diagram{\begin{tikzcd}[column sep=large]
\Omega\arrow[r,"-v_1"]&\underline{\operatorname{Hom}}(\Omega,\operatorname{Sym}^2(\Omega))\arrow[d,phantom,"\simeq"]\\
&\underline{\operatorname{Hom}}(\Omega^1_{X/S},\operatorname{Sym}^2\Omega^1_{X/S})
\end{tikzcd}}
et l'iso découlant de l'iso $\alpha$. Donc si la dim relative
